% -*- coding: utf-8 -*-

\chapter{绪论}

\section{研究背景与意义}

随着大语言模型逐步从单轮问答扩展到长时多轮交互，基于模型构建能够持续感知环境、保持用户画像并完成复杂任务的智能体系统，已经成为人工智能应用的重要发展方向。在这一过程中，长期记忆能力逐渐从可选增强模块转变为系统级基础设施。尤其在跨会话对话、长期个性化服务、复杂任务规划和多阶段协同场景下，若系统无法稳定保存并调用历史信息，就很容易出现事实遗忘、上下文断裂、人格漂移、重复提问和答非所问等问题。

从工程实现角度看，当前长期记忆系统已经形成了多种技术路线，例如基于检索增强生成的方案、基于摘要压缩的方案、基于时间层级结构组织的方案，以及引入反思与重要性评分的混合记忆方案。这些方案在系统结构与性能表现上差异明显，但它们都会面临一个共同问题：当智能体回答错误时，开发者往往只能看到结果错了，却很难进一步判断问题究竟发生在记忆写入、原生检索还是答案生成环节。若评测系统缺乏细粒度归因能力，后续优化就只能依赖经验试错，难以形成高效、可复用的系统调试闭环。

因此，构建一套面向真实长期记忆系统、具备统一协议、可解释输出和良好可扩展性的评测框架，具有明确的理论意义和工程价值。一方面，它能够帮助研究者在统一标准下比较不同记忆系统的优劣；另一方面，它可以帮助开发者快速定位问题来源，降低调试成本，提高系统迭代效率。本文即围绕这一目标展开设计与实现。

\section{国内外研究现状}

\subsection{长期记忆系统研究进展}

围绕长期记忆系统的研究，现有工作大致可以分为两条主线。第一条主线关注长上下文建模与外部记忆增强。Lewis 等提出的检索增强生成方法 [2] 为知识密集型任务提供了外部记忆调用范式；ReAct [3]、Toolformer [4]、Generative Agents [5]、Reflexion [6] 和 MemGPT [7] 等工作则进一步把推理、行动、工具使用和记忆管理纳入统一智能体框架。第二条主线更直接聚焦长期记忆系统本身，即如何围绕对话历史构建、更新和检索记忆，使系统在长时间交互中保持稳定表现。

根据对调研材料的整理，可以发现现有记忆系统虽在实现形式上差异很大，但其核心优化目标大体围绕三个问题展开：第一，如何从原始交互中提取并构建高价值记忆；第二，如何在新信息到来后进行更新、压缩或冲突消解；第三，如何在运行时高效检索并利用记忆完成回答。围绕这三个方向，近年来出现了多种代表性方法。Mem0 [8] 强调提取与更新的工具化决策，以及图增强记忆表示；TiMem [9] 强调时间层级的记忆整合；A-MEM 通过原子笔记与动态链接机制强化记忆演化；O-Mem 强调人格、情景与工作记忆的协同组织；还有工作通过递归摘要、流式覆写、即时编译式深度研究等方式提升长期记忆能力。尽管这些方法的内在数据结构不同，但它们最终都必须回答同样的问题：关键事实是否被写入、运行时是否被取回、最终答案是否正确使用了这些信息。

\subsection{长期记忆评测研究进展}

与记忆系统本身的发展相比，长期记忆评测近年也逐步从简单的黑箱问答走向更复杂的长期交互评测。LoCoMo [1] 已成为长期对话记忆研究中的代表性基准，其任务覆盖单跳问答、多跳推理、时间推理、开放域问题、对抗性问题以及事件摘要等多个维度。调研材料表明，LoCoMo 的重要性并不只在于数据长度，更在于它把“是否记住事实”“是否能进行跨片段整合”“是否能在无证据时正确拒答”等关键能力放入同一场景中进行测试。

除 LoCoMo 之外，LongMemEval [10] 明确指出现有长期助手评测与真实多会话交互不完全匹配，补充了信息提取、多会话推理、知识更新、时间推理和弃权等能力维度；MemoryBench [11] 强调记忆系统在用户反馈驱动下的持续学习效果；StoryBench [12] 用动态互动叙事场景测试长期记忆与多步推理；A-MEMGYM [13] 从交互式长期任务出发检验系统在持续决策中的记忆使用能力；HaluMem [14] 则进一步说明，仅看最终问答正确率并不足以定位问题，还需要区分错误到底发生在提取、更新还是问答阶段。

这些评测研究的共同趋势可以概括为三点：第一，评测目标从静态事实回忆扩展到了时间推理、知识更新、弃权、摘要与多跳整合等综合能力；第二，评测对象逐渐从纯大模型记忆能力扩展到真正的外部记忆系统；第三，越来越多工作开始意识到“追根溯源定位错误环节”比单一最终正确率更有价值。

\subsection{三层抽象的提出依据}

本文之所以定义编码层、检索层与生成层，并不是单纯从当前项目实现倒推出来的技术拆分，而是基于对大量文献与系统方案的归纳抽象。调研结果表明，现有记忆系统无论采用摘要式、图结构式、分层时间式、原子笔记式还是研究式架构，虽然内部模块名称各不相同，但从评测视角看都可以统一映射为三个关键环节：

\begin{enumerate}
    \item 编码层，对应系统是否将原始对话中的关键信息正确写入可持续记忆；
    \item 检索层，对应系统在运行时是否能取回与当前问题相关的有效记忆；
    \item 生成层，对应系统是否能够基于已取回的信息给出忠实、正确且在 NEG 问题下不过度编造的回答。
\end{enumerate}

这种三层抽象的意义在于，它并不依赖某一种记忆系统的内部实现，而是面向几乎所有长期记忆系统都成立的共性链路。只有先完成这种框架提取与通用抽象，评测系统才有可能在未来持续接入更多记忆系统，并在不为每个系统重写整套评测逻辑的前提下实现细粒度归因。

\subsection{现有工作的不足与本文切入点}

综合记忆系统研究与评测研究可以发现，当前仍存在四个突出问题。第一，尽管已有 benchmark 对记忆能力的覆盖越来越丰富，但很多实际评测仍然停留在最终正确率层面，难以回答错误究竟源于写入、检索还是生成。第二，现有很多分析方法与具体系统内部结构绑定较深，不利于形成统一的跨系统评测协议。第三，真实工程中的环境冲突、接口差异、中间结果不可见与长任务卡顿问题，在已有论文中通常缺少系统性讨论。第四，面向未来新系统扩展时，缺乏一种稳定的通用评测抽象，导致每增加一种系统就需要重新组织评测过程。

因此，本文将长期记忆评测定位为一个方法设计与工程系统实现并重的问题。本文的切入点不是再提出一种新的记忆模型，而是提出一个面向通用记忆系统的三层归因评测框架，并围绕统一抽象、统一协议、统一判分和统一结果组织方式展开实现。

\section{研究目标与主要内容}

本文以长期记忆智能体评测为研究对象，目标是在不改变被测系统核心逻辑的前提下，设计并实现一套可扩展、可归因、可复现的统一评测框架。具体目标包括：

\begin{enumerate}
    \item 设计统一的数据样本结构、状态空间和缺陷编码体系；
    \item 构建编码、检索、生成三个并行探针及最终归因代理；
    \item 设计统一适配器协议，并以已有真实系统作为当前接入样例验证其通用性；
    \item 建立 baseline、build、eval 分离运行方式，提高实验可复现性；
    \item 形成逐题 JSON、索引文件和运行摘要等多粒度结果组织方式；
    \item 在量化实验尚未完全补齐的情况下，为后续结果补充预留标准化位置。
\end{enumerate}

围绕上述目标，本文的主要工作可归纳为以下四点：

\begin{enumerate}
    \item 提出一套三探针并行、最终归因收敛的长期记忆评测框架；
    \item 构建统一适配器层，使异构记忆系统能够纳入同一评测协议；
    \item 在工程实现中引入统一 CorrectnessJudge、逐题结果落盘和结构化可观测性增强；
    \item 给出完整实验设计、结果记录模板与工程案例分析方法。
\end{enumerate}

\section{论文结构安排}

本文共分为八章，各章节内容安排如下：

\begin{enumerate}
    \item 第一章为绪论，介绍研究背景、研究现状、研究目标与论文结构；
    \item 第二章为需求分析与总体设计，从系统工程角度分析评测框架的目标与约束；
    \item 第三章为评测任务与指标体系设计，定义 POS/NEG 任务、状态空间和缺陷编码；
    \item 第四章为三探针归因评测框架设计，详细阐述四个代理的职责与协作关系；
    \item 第五章为多系统适配与工程实现，介绍统一适配器、脚本和结果组织机制；
    \item 第六章为实验设计与结果分析，给出实验问题、环境、指标和结果模板；
    \item 第七章为案例分析与工程讨论，分析框架在真实运行中的诊断价值；
    \item 第八章为总结与展望，总结全文工作并讨论后续研究方向。
\end{enumerate}

\chapter{需求分析与总体设计}

\section{问题场景定义}

本文所评测的对象是具备长期记忆能力的智能体系统。在这一场景中，系统需要处理跨会话、多时间跨度的历史对话，并在后续查询到来时检索并组织相关信息，最终生成与上下文一致的答案。与只评估模型最终回答文本的传统问答任务不同，本文关心的是整个长期记忆链路是否可靠，即原始信息是否被正确写入、是否在回答时被正确检索、是否被生成模块忠实利用。

从运行流程看，该问题可抽象为以下四个阶段：

\begin{enumerate}
    \item 原始多轮对话被导入记忆系统；
    \item 评测框架构建标准化评测样本；
    \item 被测系统执行在线检索与回答；
    \item 评测层从编码、检索、生成三个角度进行诊断并输出最终归因。
\end{enumerate}

这种评测问题本质上既包含算法判断，又包含运行协议、日志组织和系统适配，因此需要统一的系统设计。

\section{功能需求分析}

为支撑上述问题场景，评测框架至少应满足表 \ref{tab:function-req} 所示功能需求。

\begin{table}[htbp]
\centering
\caption{评测框架功能需求}
\label{tab:function-req}
\begin{tabular}{p{0.12\textwidth}p{0.2\textwidth}p{0.56\textwidth}}
\toprule
编号 & 需求项 & 需求说明 \\
\midrule
F1 & 数据标准化 & 将 LoCoMo 原始对话、问题、证据、答案与任务类型统一组织为标准样本对象。 \\
F2 & 多探针评测 & 支持编码、检索、生成三个并行探针分别输出状态、缺陷和证据。 \\
F3 & 最终归因 & 基于三探针结果形成统一结论，而不是只返回局部判断。 \\
F4 & 多系统适配 & 通过统一接口接入不同记忆系统，减少为单个系统重写评测逻辑的成本。 \\
F5 & 语义判分 & baseline 与 eval 共享统一语义正确性口径，避免仅依赖字符串匹配。 \\
F6 & 结果落盘 & 支持聚合结果、逐题 JSON、运行摘要与索引文件等多粒度输出。 \\
F7 & 可观测性 & 支持逐题日志、错误记录与工件引用，方便定位卡顿与适配异常。 \\
\bottomrule
\end{tabular}
\end{table}

\section{非功能需求分析}

除功能需求外，本文框架还需要满足以下非功能要求。

\subsection{可扩展性}

评测框架不能只服务于单一记忆系统，而应能够接入未来更多系统。因此，数据层、评测层与适配层必须明确分离，避免逻辑深度耦合在某个被测系统内部。

\subsection{可复现性}

真实实验常常受到外部 API、运行环境和长任务执行时间影响。为了保证结果可复现，评测框架需要显式记录配置、输出结构和中间工件。例如，对于建库开销较高的 MemBox，必须允许 build 与 eval 分离，并复用 build artifact。

\subsection{可解释性}

长期记忆系统的优化依赖对错误来源的准确把握。因此，评测结果不仅要有最终结论，还必须保留能够支撑判断的中间证据和简要归因逻辑。

\subsection{鲁棒性}

由于真实系统运行中可能出现依赖冲突、网络超时、接口异常和空结果返回等问题，评测框架必须具备错误记录、超时控制、日志心跳和降级策略，以避免整批实验因单点故障全部失败。

\section{总体架构设计}

基于以上需求，本文将评测框架划分为数据层、适配器层、评测核心层与输出层四个部分，其关系如图 \ref{fig:overall-arch} 所示。

\begin{figure}[htbp]
\centering
\fbox{\parbox{0.92\textwidth}{\centering
LoCoMo 数据集与样本构建层\\[0.6em]
$\Downarrow$\\[0.6em]
统一适配器层（O-Mem / MemBox / 未来系统）\\[0.6em]
$\Downarrow$\\[0.6em]
评测核心层：EncodingAgent、RetrievalAgent、GenerationAgent、AttributionAgent\\[0.6em]
$\Downarrow$\\[0.6em]
结果输出层：run summary、question index、per-question JSON、日志与工件引用}}
\caption{评测框架总体架构}
\label{fig:overall-arch}
\end{figure}

在该架构下，数据层负责样本构建，适配器层负责把不同系统映射到统一协议，评测核心层负责执行多探针分析与归因，输出层负责以结构化形式记录结果，最终形成从数据到结论的完整闭环。

\section{本章小结}

本章从问题定义、功能需求、非功能需求和总体架构四个方面对长期记忆智能体评测框架进行了系统分析。可以看出，该问题不仅需要方法设计，还需要围绕可扩展性、可解释性和可复现性进行系统工程建模。后续章节将在此基础上展开任务与指标设计。

\chapter{评测任务与指标体系设计}

\section{数据集与样本构成}

本文选择 LoCoMo 作为主要实验数据基础。LoCoMo 是面向超长时间跨度多轮对话记忆的评测基准 [1]，其会话长度、时间跨度和问题类型均能够较好覆盖长期记忆智能体在真实应用中会面临的核心挑战。为了让不同记忆系统都能在统一框架下接受评测，本文将原始数据统一转换为标准评测样本。每个样本至少包含以下字段：问题文本、标准答案、任务类型、关键事实键集合、原文证据、时间增强证据、oracle context 以及会话与问题标识。

这种样本设计的好处在于，后续编码、检索、生成三个探针都可以面向相同契约工作，而不再直接依赖原始数据文件的组织细节。

\section{POS 与 NEG 任务定义}

为了更准确地刻画系统行为，本文把问题划分为 POS 与 NEG 两类。

\begin{enumerate}
    \item POS 问题表示原始对话中存在支撑当前回答的证据，系统应当正确检索并作答；
    \item NEG 问题表示原始对话中不存在支撑当前回答的证据，系统理想行为应为拒答或避免编造。
\end{enumerate}

之所以必须区分这两类任务，是因为它们对系统的要求并不相同。POS 任务强调记得住且找得到，NEG 任务强调没有依据时不应凭空作答。若使用统一规则处理两类任务，会使系统在该答和不该答两个维度上的行为被混淆，削弱评测结论的解释力。

\begin{table}[htbp]
\centering
\caption{POS 与 NEG 任务差异}
\label{tab:pos-neg}
\begin{tabular}{p{0.18\textwidth}p{0.32\textwidth}p{0.32\textwidth}}
\toprule
维度 & POS 任务 & NEG 任务 \\
\midrule
原文证据 & 存在直接或可组合证据 & 不存在足以支撑回答的证据 \\
理想编码结果 & 记忆中应存在正确记录 & 不应出现可诱导回答的伪记忆 \\
理想检索结果 & 应命中关键证据 & 不应返回误导性噪声 \\
理想生成结果 & 应忠实回答 & 应谨慎拒答或不作虚构回答 \\
\bottomrule
\end{tabular}
\end{table}

\section{状态空间与缺陷编码}

为了让不同系统输出可比较、可统计，本文为三个探针定义了统一状态空间与缺陷编码。编码探针状态包括 \texttt{EXIST}、\texttt{MISS}、\texttt{CORRUPT\_AMBIG}、\texttt{CORRUPT\_WRONG} 与 \texttt{DIRTY}；检索探针状态包括 \texttt{HIT}、\texttt{MISS} 与 \texttt{NOISE}；生成探针状态包括 \texttt{PASS} 与 \texttt{FAIL}。

\begin{table}[htbp]
\centering
\caption{主要缺陷编码说明}
\label{tab:defects}
\begin{tabular}{p{0.12\textwidth}p{0.72\textwidth}}
\toprule
编码 & 含义 \\
\midrule
EM & Encoding Missing，目标记忆未正确写入。 \\
EA & Encoding Ambiguous，记忆存在但主体、指代或边界模糊。 \\
EW & Encoding Wrong，记忆存在但关键事实值错误。 \\
DMP & Dirty Memory Pollution，在 NEG 任务下检测到可诱导错误回答的伪记忆。 \\
RF & Retrieval Failure，在编码层非缺失前提下检索未命中目标记忆。 \\
LATE & 目标记忆命中但排序过晚。 \\
NOI & 检索结果噪声过多，影响系统使用。 \\
NIR & Negative Irrelevant Retrieval，在 NEG 任务下返回误导性噪声。 \\
GH & Generation Hallucination，NEG 任务下本应拒答但产生不可信回答。 \\
GF & Generation Faithfulness Failure，POS 任务下未忠实利用证据。 \\
GRF & Generation Reasoning Failure，POS 任务下读取到证据但推理或整合失败。 \\
\bottomrule
\end{tabular}
\end{table}

通过这组统一编码，系统错误不再被压缩成单一答错，而是能够映射到更明确的责任位置。

\section{统一语义判分机制}

在长期对话场景中，仅用字符串完全匹配判断答案是否正确存在明显局限。例如，同义表达、时间格式改写、简称与全称转换都可能导致正确答案被误判为错误。为此，本文在实现中引入统一 CorrectnessJudge，将规则判分与大模型语义判分结合起来，用于 baseline 与 eval 的统一口径。考虑到当前主流 API 服务通常建立在 GPT-4 类模型能力之上 [15]，本文在评测工程中把语义判分模块设计为与具体后端解耦的可配置组件。TiMem 对语义判定提示词的组织方式 [9] 为本文设计语义判分提示提供了重要参考。

设评测样本集合为 $Q=\{q_1,q_2,\ldots,q_N\}$，则最终正确率定义为：

\begin{equation}
Acc_{final}=\frac{1}{N}\sum_{i=1}^{N}\mathbb{I}(final\_correct(q_i)=1)
\end{equation}

其中，$final\_correct(q_i)$ 表示第 $i$ 个问题在统一语义判分后的最终正确性。除总体正确率外，本文还统计缺陷分布。设缺陷 $d$ 的出现频次为：

\begin{equation}
Freq(d)=\sum_{i=1}^{N}\mathbb{I}(d \in defects(q_i))
\end{equation}

该指标直接反映系统主要瓶颈是来自编码、检索还是生成环节。

\section{最终归因规则}

最终归因代理的任务不是重复输出三个探针已经给出的详细证据，而是在已有结果基础上形成稳定、简洁且面向开发者可读的结论。本文将最终归因收敛为以下字段：\texttt{final\_attribution}、\texttt{primary\_cause}、\texttt{secondary\_causes}、\texttt{decision\_logic} 和 \texttt{final\_judgement}。

在规则层面，本文遵循先编码、后检索、再生成的责任传导原则。例如，当编码层已经为 \texttt{MISS} 时，检索层的 \texttt{RF} 不再被视为主要故障；只有编码层非缺失而检索层未命中时，\texttt{RF} 才成立。该机制可以有效避免错误归因层层叠加。

\section{本章小结}

本章围绕数据样本、任务类型、状态空间、缺陷编码和统一语义判分完成了指标体系设计。通过对 POS/NEG 任务的显式区分，以及对三层失效的结构化描述，本文为后续的框架设计与工程实现奠定了形式化基础。

\chapter{三探针归因评测框架设计}

\section{设计原则}

为了在真实系统上获得可靠、可解释的评测结果，本文提出以下设计原则：

\begin{enumerate}
    \item 保真性原则：尽量复用被测系统原生检索与原生生成逻辑，避免评测器替代系统行为；
    \item 分层性原则：编码、检索、生成各自负责局部判断，最终归因负责全局收敛；
    \item 并行性原则：三个探针尽可能并行执行，以降低评测延迟并保持状态独立性；
    \item 可追踪性原则：每一层输出都保留结构化证据，便于后续人工复核。
\end{enumerate}

\section{编码探针设计}

编码探针用于回答原始对话中的关键信息是否被系统正确写入记忆库。对 POS 任务，它需要判断目标事实是否存在、是否模糊以及是否写错；对 NEG 任务，它需要检查系统是否误写入了本不该存在的伪记忆。

在实现上，编码探针会综合以下几类观测来源：适配器导出的全量记忆视图、适配器原生候选结果、框架侧混合候选、原生检索阴影集以及外部高召回检索器返回的补充候选。外部高召回接口的引入使本文框架能够在不改动评测核心逻辑的情况下接入后续自定义 RAG 检索方案，提高编码观测完整性。

\section{检索探针设计}

检索探针关注的问题是系统在在线回答时是否真正取回了应当使用的记忆证据。与编码探针不同，检索探针关注的是运行时原生检索输出，而非记忆库是否存在目标记录。

对于 POS 任务，若关键证据被成功检出，则状态为 \texttt{HIT}；若未命中，则状态为 \texttt{MISS}；若虽然命中但存在大量误导性噪声，则可视为 \texttt{NOISE}。对于 NEG 任务，检索探针重点观察系统是否返回足以诱导回答的无关信息，此时可归为 \texttt{NIR}。该设计使检索失败与编码失败可以被明确区分。

\section{生成探针设计}

生成探针关注的是系统在给定上下文条件下是否生成了正确且忠实的答案。为了分离检索误差和生成误差，本文同时比较在线回答与 oracle 回答。在线回答来自真实系统运行路径，oracle 回答则在提供理想证据上下文条件下生成。

对于 POS 任务，若 oracle 回答本身就错误，则说明生成阶段存在较明显缺陷；若在线回答相较 oracle 回答进一步偏离，则可区分为未忠实利用证据的 \texttt{GF} 与读取到证据但推理失败的 \texttt{GRF}。对于 NEG 任务，若系统在无依据条件下仍给出肯定式回答，则可归为 \texttt{GH}。

\section{最终归因代理设计}

最终归因代理接收三个探针输出，并按照规则进行冲突消解与因果收敛。其主要职责包括：第一，对局部结果进行冲突处理，例如在编码层缺失时抑制 \texttt{RF}；第二，把复杂中间证据压缩为开发者可直接理解的最终判断。

\begin{table}[htbp]
\centering
\caption{最终归因代理输出字段}
\label{tab:attr-output}
\begin{tabular}{p{0.22\textwidth}p{0.58\textwidth}}
\toprule
字段名 & 说明 \\
\midrule
\texttt{final\_attribution} & 最终归因标签，用于快速查看主要失效环节。 \\
\texttt{primary\_cause} & 主因，例如 encoding、retrieval、generation 或 no\_major\_failure。 \\
\texttt{secondary\_causes} & 次要原因列表，用于补充复杂情形。 \\
\texttt{decision\_logic} & 面向人工复核的简要判断逻辑。 \\
\texttt{final\_judgement} & 对系统回答与 oracle 回答的综合结论。 \\
\bottomrule
\end{tabular}
\end{table}

\section{评测流程设计}

综合三探针与归因代理，本文评测主流程可概括为以下步骤：

\begin{enumerate}
    \item 根据样本所属会话准备或复用系统运行上下文；
    \item 并行执行编码探针、检索探针和生成探针；
    \item 使用统一语义判分模块计算最终正确性；
    \item 调用最终归因代理融合三个探针的状态与缺陷；
    \item 将逐题结果写入结构化 JSON 文件；
    \item 汇总运行摘要、错误索引和缺陷统计。
\end{enumerate}

上述流程既保持了三层诊断的独立性，又能通过统一输出结构把结果收敛为面向开发者可读的评测报告。

\section{本章小结}

本章给出了三探针并行评测框架与最终归因代理的设计细节。该设计既保留了不同阶段的独立判断能力，又通过统一归因代理形成了对系统错误的整体解释，是本文方法部分的核心。

\chapter{多系统适配与工程实现}

\section{项目总体实现结构}

本文实现的项目由数据层、评测核心层、适配器层、流水线层、脚本层和测试层构成。数据层负责构建 LoCoMo 标准样本；评测核心层负责多探针分析与归因；适配器层负责对接真实记忆系统；流水线层负责组织端到端运行；脚本层提供统一 CLI 入口；测试层用于验证关键行为。项目总体说明与运行入口在仓库说明文档中已有整理 [16][17]。

\begin{table}[htbp]
\centering
\caption{项目主要模块与职责}
\label{tab:modules}
\begin{tabular}{p{0.2\textwidth}p{0.58\textwidth}}
\toprule
模块 & 职责 \\
\midrule
\texttt{dataset} & 构建 LoCoMo 标准评测样本。 \\
\texttt{eval\_core} & 实现三探针、归因代理、语义判分与高召回扩展。 \\
\texttt{adapters} & 为 O-Mem、MemBox 等系统提供统一适配。 \\
\texttt{pipeline} & 串联样本、适配器、评测核心与结果落盘。 \\
\texttt{scripts} & 提供 baseline、build、eval 等命令行入口。 \\
\texttt{tests} & 校验核心状态转移、归因规则和适配器注册。 \\
\bottomrule
\end{tabular}
\end{table}

\section{统一适配器协议设计}

为了使评测框架能够接入不同记忆系统，本文设计了统一适配器协议。一个适配器至少需要支持以下能力：导入标准化对话并建立运行态、导出全量记忆视图、查找与问题相关的记忆候选、执行原生检索、执行在线回答、执行 oracle context 下的理想回答，以及在需要时导出或加载 build 工件。

该协议一方面保持了对不同系统内部实现的兼容性，另一方面保证评测层的逻辑可以尽量保持统一，从而降低多系统接入的额外开发成本。

\section{通用适配类型一：在线检索与在线回答型系统}

一类长期记忆系统在运行时更强调原生在线检索与原生回答能力。对于这类系统，适配器的重点不是重建其中间逻辑，而是尽量保留原系统的写入、检索与回答路径，把评测层放在系统外部。此类适配通常需要解决三类问题：其一是运行环境和依赖冲突，例如深度学习组件版本不匹配导致系统无法启动；其二是配置显式化，即把模型路径、API 凭据和运行参数从系统内部抽离出来；其三是结果标准化，即把原生输出统一映射到评测层可理解的证据结构。

本文当前接入的部分真实系统可以归入这一类。它们说明：当评测目标是尽可能适配更多系统时，最重要的不是为某个系统单独定制复杂逻辑，而是抽象出一组通用适配原则，使外部评测层能够稳定读取原系统行为。

\section{通用适配类型二：建库与问答分阶段系统}

另一类长期记忆系统具有较重的离线建库过程，系统在正式问答前会先生成中间索引、记忆文件或阶段性工件。对于这类系统，若仍采用“每次评测都重新建库”的方式，不仅重复开销较高，而且很容易把建库耗时误判为在线评测卡顿。因此，本文将这类系统统一抽象为“build 与 eval 分阶段”模式：先显式完成建库，再复用 build artifact 进行 baseline 或 eval。

\begin{figure}[htbp]
\centering
\fbox{\parbox{0.88\textwidth}{\centering
原始对话导入\\[0.4em]
$\Downarrow$\\[0.4em]
\texttt{build}：构建记忆库并导出 manifest\\[0.4em]
$\Downarrow$\\[0.4em]
\texttt{baseline}/\texttt{eval}：加载 manifest，复用 build artifact 进行问答与评测}}
\caption{MemBox 的 build/eval 分离流程}
\label{fig:membox-build}
\end{figure}

这一设计不仅降低了重复建库成本，还增强了实验的阶段可观测性，使开发者更容易区分问题究竟发生在 build 还是 eval 环节。本文当前接入的部分系统可以视为这一适配类型的实例，但论文在方法表述上并不依赖某一特定系统，而是强调对“分阶段运行型记忆系统”的通用支持。

\section{结果组织与可观测性实现}

为了提高评测框架的实际可用性，本文在输出层实现了多粒度结果落盘：聚合结果文件用于全局汇总，\texttt{run\_summary.json} 用于记录运行摘要，\texttt{question\_index.json} 用于建立逐题索引，每道题还会单独生成 JSON 文件，保存系统答案、oracle 答案、探针状态、缺陷编码、最终归因和工件引用。

同时，评测脚本在 \texttt{build}、\texttt{baseline} 和 \texttt{eval} 过程中输出结构化心跳日志，使用户能够直接看到当前处理进度、系统回答和正确性判断。这一机制对长任务调试和异常排查具有重要价值。

\section{测试与验证机制}

除端到端运行外，本文还围绕关键逻辑构建了自动化测试，包括编码、检索、生成与归因代理的规则测试，\texttt{RF} 抑制与补加测试，外部高召回检索器接入测试，适配器注册与 manifest 导出测试，以及 mock pipeline 冒烟测试。测试本身不能替代真实实验，但能够保证框架迭代时不破坏核心行为。

\section{本章小结}

本章从统一适配器协议、O-Mem 与 MemBox 的接入方式、结果组织和测试机制四个方面介绍了框架实现细节。通过这些实现，本文把评测方法真正转化为可运行、可复现、可扩展的工程系统。

\chapter{实验设计与结果分析}

\section{本章说明}

由于当前实验尚未最终完成，本文本章暂不展开正式实验结果、结果分析与统计讨论，以避免在论文中提前写入未经最终核验的数据。后续在完成 O-Mem、MemBox 及更多记忆系统的统一跑测后，将补写以下内容：

\begin{enumerate}
    \item 实验环境与数据配置；
    \item 对比方案与评价指标；
    \item 总体结果表、缺陷分布表与案例分析；
    \item 对系统差异、框架有效性与通用性的实验讨论。
\end{enumerate}

\section{待补内容位置}

为保证后续补写结构统一，本章预留如下补写位置：

\begin{enumerate}
    \item 长期记忆系统总体结果对比；
    \item 三层探针的状态与缺陷分布；
    \item 通用适配框架在不同系统上的运行稳定性；
    \item 典型问题的逐题归因案例。
\end{enumerate}

\section{本章小结}

本章暂作留白处理，待实验结果全部整理完成后补入正式内容。

\chapter{案例分析与工程讨论}

\section{案例一：在线检索型系统的环境鲁棒性问题}

对于强调原生在线检索和原生回答的记忆系统，最常见的工程问题并不一定发生在评测逻辑内部，而可能发生在系统启动阶段。例如，嵌入模型加载常常受到 \texttt{flash-attn}、\texttt{torch} 与 \texttt{transformers} 等依赖兼容关系影响。若缺少统一适配层与结构化日志，这类问题往往会被误判为评测框架不可用。本文当前接入的部分真实系统就呈现出这一特征。

本文实现的统一评测框架并不掩盖这些真实运行问题，而是通过适配器层、运行脚本和错误记录机制，把环境问题、运行问题与评测问题分开记录。这样，开发者能够区分系统未能正常启动和系统启动后回答错误两类问题，从而更高效地组织调试流程。

\section{案例二：分阶段系统的可观测性问题}

对于需要先建库、再问答的长期记忆系统，一个典型工程问题是前置阶段耗时较长。若直接把 build 和 eval 混合在同一个长任务中，用户往往只能看到任务长时间没有输出，难以判断当前处于正常建库、异常卡住还是在线评测超时状态。本文通过 build artifact 的引入，把这一问题显式拆分。当前已接入的部分真实系统正好体现了这一类问题。

在该设计下，用户可以先确认 build 是否成功生成 manifest，再复用该工件运行 baseline 和 eval。若 eval 阶段仍然耗时异常，则可以进一步通过逐题心跳日志和逐题结果文件判断问题发生在具体哪一道题。由此，系统问题被从整体不可见转化为阶段可区分、逐题可定位的可观察问题。

\section{多探针归因的诊断价值}

与仅输出正确率相比，多探针归因的价值主要体现在以下三个方面：

\begin{enumerate}
    \item 当编码层大量出现 \texttt{EM} 或 \texttt{EW} 时，说明应优先优化记忆写入与压缩策略；
    \item 当检索层出现较多 \texttt{RF}、\texttt{LATE} 或 \texttt{NOI} 时，说明应优先优化检索召回与排序；
    \item 当生成层以 \texttt{GF} 或 \texttt{GRF} 为主时，说明问题更多位于答案忠实性或推理整合能力，而非记忆库本身。
\end{enumerate}

因此，三探针归因不仅是一种评测方法，也是一种面向系统调优的诊断工具。

\section{适用范围与局限}

尽管本文框架已经能够覆盖多种类型的真实记忆系统，但仍存在若干局限：

\begin{enumerate}
    \item 当前主实验对象主要基于 LoCoMo，尚未扩展到更多真实用户数据；
    \item 归因规则虽然结构化清晰，但对极端复杂失败链仍可能存在简化；
    \item 统一语义判分提升了鲁棒性，但也引入了评审模型本身的偏差；
    \item 某些系统内部中间状态无法完全导出，因此评测仍然存在观测盲区。
\end{enumerate}

这些局限不会削弱本文作为工程化评测框架的主体价值，但的确为后续研究留下了进一步拓展空间。

\section{本章小结}

本章结合两类典型工程案例，讨论了统一评测框架在真实开发和问题定位中的价值。可以看出，本文方法不仅服务于最终实验统计，更能在系统开发与调试阶段持续发挥作用。

\chapter{总结与展望}

\section{本章说明}

由于实验结果与案例统计仍待最终补齐，本文总结与展望部分暂不写入完整结论性表述，以避免在最终数据尚未稳定之前形成过早归纳。待实验部分补齐后，本章将围绕以下内容进行正式撰写：

\begin{enumerate}
    \item 本文工作的总体贡献总结；
    \item 三层评测框架的有效性与通用性归纳；
    \item 当前方法的局限性说明；
    \item 面向未来新记忆系统与新 benchmark 的扩展展望。
\end{enumerate}

\section{待补说明}

本章当前留白，后续将根据最终实验结果和整体论文定稿情况补写。
